% =============================================================================
% Captures d'écran — dossier mecano_cap (Overleaf : uploadez le dossier à la racine
% du projet, ou copiez les PNG dans un sous-dossier "figures" et adaptez le chemin).
% =============================================================================
% Dans le préambule (global_config.tex ou main.tex), assurez-vous d'avoir :
%   \usepackage{graphicx}
%   \usepackage{caption}      % optionnel : sous-titres plus propres
%   \usepackage{subcaption}     % optionnel : côte à côte
%
% Chemin des images : choisissez UNE des deux lignes suivantes (décommentez).
% =============================================================================

% Si vous changez le dossier, modifiez la ligne suivante (ex. {figures/}).
\graphicspath{{mecano_cap/}}

% Largeur par défaut des captures (ajustez 0.85 si trop grand).
\newcommand{\CapWidth}{0.92\textwidth}

% Usage : \Screenshot{label}{fichier_sans_chemin}{légende courte}
\newcommand{\Screenshot}[3]{%
  \begin{figure}[htbp]
    \centering
    \includegraphics[width=\CapWidth]{#2}%
    \caption{#3}%
    \label{fig:#1}%
  \end{figure}%
}

% Fichiers avec ESPACES dans le nom : passer le nom exact entre accolades.
\newcommand{\fichierListeClient}{{liste client.png}}
\newcommand{\fichierListeStock}{{liste stock.png}}

% -----------------------------------------------------------------------------
% CHAPITRE — Introduction / général (connexion, tableau de bord)
% -----------------------------------------------------------------------------
% % =============================================================================
% Captures d'écran — dossier mecano_cap (Overleaf : uploadez le dossier à la racine
% du projet, ou copiez les PNG dans un sous-dossier "figures" et adaptez le chemin).
% =============================================================================
% Dans le préambule (global_config.tex ou main.tex), assurez-vous d'avoir :
%   \usepackage{graphicx}
%   \usepackage{caption}      % optionnel : sous-titres plus propres
%   \usepackage{subcaption}     % optionnel : côte à côte
%
% Chemin des images : choisissez UNE des deux lignes suivantes (décommentez).
% =============================================================================

% Si vous changez le dossier, modifiez la ligne suivante (ex. {figures/}).
\graphicspath{{mecano_cap/}}

% Largeur par défaut des captures (ajustez 0.85 si trop grand).
\newcommand{\CapWidth}{0.92\textwidth}

% Usage : \Screenshot{label}{fichier_sans_chemin}{légende courte}
\newcommand{\Screenshot}[3]{%
  \begin{figure}[htbp]
    \centering
    \includegraphics[width=\CapWidth]{#2}%
    \caption{#3}%
    \label{fig:#1}%
  \end{figure}%
}

% Fichiers avec ESPACES dans le nom : passer le nom exact entre accolades.
\newcommand{\fichierListeClient}{{liste client.png}}
\newcommand{\fichierListeStock}{{liste stock.png}}

% -----------------------------------------------------------------------------
% CHAPITRE — Introduction / général (connexion, tableau de bord)
% -----------------------------------------------------------------------------
% % =============================================================================
% Captures d'écran — dossier mecano_cap (Overleaf : uploadez le dossier à la racine
% du projet, ou copiez les PNG dans un sous-dossier "figures" et adaptez le chemin).
% =============================================================================
% Dans le préambule (global_config.tex ou main.tex), assurez-vous d'avoir :
%   \usepackage{graphicx}
%   \usepackage{caption}      % optionnel : sous-titres plus propres
%   \usepackage{subcaption}     % optionnel : côte à côte
%
% Chemin des images : choisissez UNE des deux lignes suivantes (décommentez).
% =============================================================================

% Si vous changez le dossier, modifiez la ligne suivante (ex. {figures/}).
\graphicspath{{mecano_cap/}}

% Largeur par défaut des captures (ajustez 0.85 si trop grand).
\newcommand{\CapWidth}{0.92\textwidth}

% Usage : \Screenshot{label}{fichier_sans_chemin}{légende courte}
\newcommand{\Screenshot}[3]{%
  \begin{figure}[htbp]
    \centering
    \includegraphics[width=\CapWidth]{#2}%
    \caption{#3}%
    \label{fig:#1}%
  \end{figure}%
}

% Fichiers avec ESPACES dans le nom : passer le nom exact entre accolades.
\newcommand{\fichierListeClient}{{liste client.png}}
\newcommand{\fichierListeStock}{{liste stock.png}}

% -----------------------------------------------------------------------------
% CHAPITRE — Introduction / général (connexion, tableau de bord)
% -----------------------------------------------------------------------------
% % =============================================================================
% Captures d'écran — dossier mecano_cap (Overleaf : uploadez le dossier à la racine
% du projet, ou copiez les PNG dans un sous-dossier "figures" et adaptez le chemin).
% =============================================================================
% Dans le préambule (global_config.tex ou main.tex), assurez-vous d'avoir :
%   \usepackage{graphicx}
%   \usepackage{caption}      % optionnel : sous-titres plus propres
%   \usepackage{subcaption}     % optionnel : côte à côte
%
% Chemin des images : choisissez UNE des deux lignes suivantes (décommentez).
% =============================================================================

% Si vous changez le dossier, modifiez la ligne suivante (ex. {figures/}).
\graphicspath{{mecano_cap/}}

% Largeur par défaut des captures (ajustez 0.85 si trop grand).
\newcommand{\CapWidth}{0.92\textwidth}

% Usage : \Screenshot{label}{fichier_sans_chemin}{légende courte}
\newcommand{\Screenshot}[3]{%
  \begin{figure}[htbp]
    \centering
    \includegraphics[width=\CapWidth]{#2}%
    \caption{#3}%
    \label{fig:#1}%
  \end{figure}%
}

% Fichiers avec ESPACES dans le nom : passer le nom exact entre accolades.
\newcommand{\fichierListeClient}{{liste client.png}}
\newcommand{\fichierListeStock}{{liste stock.png}}

% -----------------------------------------------------------------------------
% CHAPITRE — Introduction / général (connexion, tableau de bord)
% -----------------------------------------------------------------------------
% \input{mecano_cap_figures} puis décommentez les blocs voulus dans chaque chapitre.

\newcommand{\FiguresIntroGenerale}{%
  \Screenshot{login}{login}{Écran de connexion à l'application}%
  \Screenshot{dashboard}{dashboard}{Tableau de bord principal}%
}

% -----------------------------------------------------------------------------
% CHAPITRE 1 — Contexte, cadre, solution (vue d'ensemble métier)
% -----------------------------------------------------------------------------
\newcommand{\FiguresChapitreUn}{%
  \Screenshot{dashboard_ctx}{dashboard}{Vue synthétique de l'activité (illustration du contexte numérique)}%
}

% -----------------------------------------------------------------------------
% CHAPITRE 2 — Acteurs, exigences (administration / utilisateurs)
% -----------------------------------------------------------------------------
\newcommand{\FiguresChapitreDeux}{%
  \Screenshot{users_liste}{user}{Gestion des utilisateurs et profils}%
  \Screenshot{ajout_user}{ajout_user}{Formulaire d'ajout d'utilisateur}%
}

% -----------------------------------------------------------------------------
% CHAPITRE 3 — Release 1 : véhicules, archives, ordre de réparation (atelier)
% -----------------------------------------------------------------------------
\newcommand{\FiguresChapitreTrois}{%
  \Screenshot{liste_vehicules}{liste_vehicule}{Liste des véhicules}%
  \Screenshot{ajout_vehicule}{ajout_vehicule}{Ajout d'un véhicule}%
  \Screenshot{changement_etat}{changement_etat}{Changement d'état / suivi du véhicule}%
  \Screenshot{suivre_vehicule}{suivre_vehicule}{Fiche ou suivi associé au véhicule}%
  \Screenshot{affectation}{affectation}{Affectation (planning / équipe)}%
}

% -----------------------------------------------------------------------------
% CHAPITRE 4 — Release 2 : clients, devis, stock, facturation, réclamations, calendrier
% -----------------------------------------------------------------------------
\newcommand{\FiguresChapitreQuatre}{%
  \Screenshot{nv_client}{nv_client}{Création / fiche client}%
  \Screenshot{liste_client}{\fichierListeClient}{Liste des clients}%
  \Screenshot{client_dettes}{client_dettes}{Clients et dettes}%
  \Screenshot{demande_devis}{demande_devis}{Demandes de devis}%
  \Screenshot{ajout_demande_devis}{ajout_demande_devis}{Ajout d'une demande de devis}%
  \Screenshot{liste_stock}{\fichierListeStock}{Liste du stock général}%
  \Screenshot{ajout_produit}{ajout_produit}{Ajout de produit}%
  \Screenshot{modif_produit}{modif_produit}{Modification de produit}%
  \Screenshot{supprimer_produit}{supprimer_produit}{Suppression de produit}%
  \Screenshot{facturation_vente}{facturation_vente}{Facturation vente}%
  \Screenshot{nouvelle_facture}{nouvelle_facture}{Création d'une nouvelle facture}%
  \Screenshot{facture_achat}{facture_achat}{Facturation / suivi achat}%
  \Screenshot{fournisseurs}{fourniseurs}{Gestion des fournisseurs (nom de fichier : fourniseurs.png)}%
  \Screenshot{reclamations}{reclamations}{Liste des réclamations}%
  \Screenshot{ajout_reclamation}{ajout_reclamation}{Ajout d'une réclamation}%
  \Screenshot{calendrier}{calendrier}{Calendrier / planning}%
}

% -----------------------------------------------------------------------------
% CHAPITRE 5 — Tests & validation (illustrations des écrans validés)
% Reprendre une sélection représentative (éviter 20 figures d'un coup).
% -----------------------------------------------------------------------------
\newcommand{\FiguresChapitreCinqTests}{%
  \Screenshot{login_test}{login}{Vérification du parcours d'authentification}%
  \Screenshot{dashboard_test}{dashboard}{Contrôle du tableau de bord après connexion}%
  \Screenshot{liste_vehicules_test}{liste_vehicule}{Scénario : navigation liste véhicules}%
  \Screenshot{facturation_test}{facturation_vente}{Scénario : module facturation}%
}
 puis décommentez les blocs voulus dans chaque chapitre.

\newcommand{\FiguresIntroGenerale}{%
  \Screenshot{login}{login}{Écran de connexion à l'application}%
  \Screenshot{dashboard}{dashboard}{Tableau de bord principal}%
}

% -----------------------------------------------------------------------------
% CHAPITRE 1 — Contexte, cadre, solution (vue d'ensemble métier)
% -----------------------------------------------------------------------------
\newcommand{\FiguresChapitreUn}{%
  \Screenshot{dashboard_ctx}{dashboard}{Vue synthétique de l'activité (illustration du contexte numérique)}%
}

% -----------------------------------------------------------------------------
% CHAPITRE 2 — Acteurs, exigences (administration / utilisateurs)
% -----------------------------------------------------------------------------
\newcommand{\FiguresChapitreDeux}{%
  \Screenshot{users_liste}{user}{Gestion des utilisateurs et profils}%
  \Screenshot{ajout_user}{ajout_user}{Formulaire d'ajout d'utilisateur}%
}

% -----------------------------------------------------------------------------
% CHAPITRE 3 — Release 1 : véhicules, archives, ordre de réparation (atelier)
% -----------------------------------------------------------------------------
\newcommand{\FiguresChapitreTrois}{%
  \Screenshot{liste_vehicules}{liste_vehicule}{Liste des véhicules}%
  \Screenshot{ajout_vehicule}{ajout_vehicule}{Ajout d'un véhicule}%
  \Screenshot{changement_etat}{changement_etat}{Changement d'état / suivi du véhicule}%
  \Screenshot{suivre_vehicule}{suivre_vehicule}{Fiche ou suivi associé au véhicule}%
  \Screenshot{affectation}{affectation}{Affectation (planning / équipe)}%
}

% -----------------------------------------------------------------------------
% CHAPITRE 4 — Release 2 : clients, devis, stock, facturation, réclamations, calendrier
% -----------------------------------------------------------------------------
\newcommand{\FiguresChapitreQuatre}{%
  \Screenshot{nv_client}{nv_client}{Création / fiche client}%
  \Screenshot{liste_client}{\fichierListeClient}{Liste des clients}%
  \Screenshot{client_dettes}{client_dettes}{Clients et dettes}%
  \Screenshot{demande_devis}{demande_devis}{Demandes de devis}%
  \Screenshot{ajout_demande_devis}{ajout_demande_devis}{Ajout d'une demande de devis}%
  \Screenshot{liste_stock}{\fichierListeStock}{Liste du stock général}%
  \Screenshot{ajout_produit}{ajout_produit}{Ajout de produit}%
  \Screenshot{modif_produit}{modif_produit}{Modification de produit}%
  \Screenshot{supprimer_produit}{supprimer_produit}{Suppression de produit}%
  \Screenshot{facturation_vente}{facturation_vente}{Facturation vente}%
  \Screenshot{nouvelle_facture}{nouvelle_facture}{Création d'une nouvelle facture}%
  \Screenshot{facture_achat}{facture_achat}{Facturation / suivi achat}%
  \Screenshot{fournisseurs}{fourniseurs}{Gestion des fournisseurs (nom de fichier : fourniseurs.png)}%
  \Screenshot{reclamations}{reclamations}{Liste des réclamations}%
  \Screenshot{ajout_reclamation}{ajout_reclamation}{Ajout d'une réclamation}%
  \Screenshot{calendrier}{calendrier}{Calendrier / planning}%
}

% -----------------------------------------------------------------------------
% CHAPITRE 5 — Tests & validation (illustrations des écrans validés)
% Reprendre une sélection représentative (éviter 20 figures d'un coup).
% -----------------------------------------------------------------------------
\newcommand{\FiguresChapitreCinqTests}{%
  \Screenshot{login_test}{login}{Vérification du parcours d'authentification}%
  \Screenshot{dashboard_test}{dashboard}{Contrôle du tableau de bord après connexion}%
  \Screenshot{liste_vehicules_test}{liste_vehicule}{Scénario : navigation liste véhicules}%
  \Screenshot{facturation_test}{facturation_vente}{Scénario : module facturation}%
}
 puis décommentez les blocs voulus dans chaque chapitre.

\newcommand{\FiguresIntroGenerale}{%
  \Screenshot{login}{login}{Écran de connexion à l'application}%
  \Screenshot{dashboard}{dashboard}{Tableau de bord principal}%
}

% -----------------------------------------------------------------------------
% CHAPITRE 1 — Contexte, cadre, solution (vue d'ensemble métier)
% -----------------------------------------------------------------------------
\newcommand{\FiguresChapitreUn}{%
  \Screenshot{dashboard_ctx}{dashboard}{Vue synthétique de l'activité (illustration du contexte numérique)}%
}

% -----------------------------------------------------------------------------
% CHAPITRE 2 — Acteurs, exigences (administration / utilisateurs)
% -----------------------------------------------------------------------------
\newcommand{\FiguresChapitreDeux}{%
  \Screenshot{users_liste}{user}{Gestion des utilisateurs et profils}%
  \Screenshot{ajout_user}{ajout_user}{Formulaire d'ajout d'utilisateur}%
}

% -----------------------------------------------------------------------------
% CHAPITRE 3 — Release 1 : véhicules, archives, ordre de réparation (atelier)
% -----------------------------------------------------------------------------
\newcommand{\FiguresChapitreTrois}{%
  \Screenshot{liste_vehicules}{liste_vehicule}{Liste des véhicules}%
  \Screenshot{ajout_vehicule}{ajout_vehicule}{Ajout d'un véhicule}%
  \Screenshot{changement_etat}{changement_etat}{Changement d'état / suivi du véhicule}%
  \Screenshot{suivre_vehicule}{suivre_vehicule}{Fiche ou suivi associé au véhicule}%
  \Screenshot{affectation}{affectation}{Affectation (planning / équipe)}%
}

% -----------------------------------------------------------------------------
% CHAPITRE 4 — Release 2 : clients, devis, stock, facturation, réclamations, calendrier
% -----------------------------------------------------------------------------
\newcommand{\FiguresChapitreQuatre}{%
  \Screenshot{nv_client}{nv_client}{Création / fiche client}%
  \Screenshot{liste_client}{\fichierListeClient}{Liste des clients}%
  \Screenshot{client_dettes}{client_dettes}{Clients et dettes}%
  \Screenshot{demande_devis}{demande_devis}{Demandes de devis}%
  \Screenshot{ajout_demande_devis}{ajout_demande_devis}{Ajout d'une demande de devis}%
  \Screenshot{liste_stock}{\fichierListeStock}{Liste du stock général}%
  \Screenshot{ajout_produit}{ajout_produit}{Ajout de produit}%
  \Screenshot{modif_produit}{modif_produit}{Modification de produit}%
  \Screenshot{supprimer_produit}{supprimer_produit}{Suppression de produit}%
  \Screenshot{facturation_vente}{facturation_vente}{Facturation vente}%
  \Screenshot{nouvelle_facture}{nouvelle_facture}{Création d'une nouvelle facture}%
  \Screenshot{facture_achat}{facture_achat}{Facturation / suivi achat}%
  \Screenshot{fournisseurs}{fourniseurs}{Gestion des fournisseurs (nom de fichier : fourniseurs.png)}%
  \Screenshot{reclamations}{reclamations}{Liste des réclamations}%
  \Screenshot{ajout_reclamation}{ajout_reclamation}{Ajout d'une réclamation}%
  \Screenshot{calendrier}{calendrier}{Calendrier / planning}%
}

% -----------------------------------------------------------------------------
% CHAPITRE 5 — Tests & validation (illustrations des écrans validés)
% Reprendre une sélection représentative (éviter 20 figures d'un coup).
% -----------------------------------------------------------------------------
\newcommand{\FiguresChapitreCinqTests}{%
  \Screenshot{login_test}{login}{Vérification du parcours d'authentification}%
  \Screenshot{dashboard_test}{dashboard}{Contrôle du tableau de bord après connexion}%
  \Screenshot{liste_vehicules_test}{liste_vehicule}{Scénario : navigation liste véhicules}%
  \Screenshot{facturation_test}{facturation_vente}{Scénario : module facturation}%
}
 puis décommentez les blocs voulus dans chaque chapitre.

\newcommand{\FiguresIntroGenerale}{%
  \Screenshot{login}{login}{Écran de connexion à l'application}%
  \Screenshot{dashboard}{dashboard}{Tableau de bord principal}%
}

% -----------------------------------------------------------------------------
% CHAPITRE 1 — Contexte, cadre, solution (vue d'ensemble métier)
% -----------------------------------------------------------------------------
\newcommand{\FiguresChapitreUn}{%
  \Screenshot{dashboard_ctx}{dashboard}{Vue synthétique de l'activité (illustration du contexte numérique)}%
}

% -----------------------------------------------------------------------------
% CHAPITRE 2 — Acteurs, exigences (administration / utilisateurs)
% -----------------------------------------------------------------------------
\newcommand{\FiguresChapitreDeux}{%
  \Screenshot{users_liste}{user}{Gestion des utilisateurs et profils}%
  \Screenshot{ajout_user}{ajout_user}{Formulaire d'ajout d'utilisateur}%
}

% -----------------------------------------------------------------------------
% CHAPITRE 3 — Release 1 : véhicules, archives, ordre de réparation (atelier)
% -----------------------------------------------------------------------------
\newcommand{\FiguresChapitreTrois}{%
  \Screenshot{liste_vehicules}{liste_vehicule}{Liste des véhicules}%
  \Screenshot{ajout_vehicule}{ajout_vehicule}{Ajout d'un véhicule}%
  \Screenshot{changement_etat}{changement_etat}{Changement d'état / suivi du véhicule}%
  \Screenshot{suivre_vehicule}{suivre_vehicule}{Fiche ou suivi associé au véhicule}%
  \Screenshot{affectation}{affectation}{Affectation (planning / équipe)}%
}

% -----------------------------------------------------------------------------
% CHAPITRE 4 — Release 2 : clients, devis, stock, facturation, réclamations, calendrier
% -----------------------------------------------------------------------------
\newcommand{\FiguresChapitreQuatre}{%
  \Screenshot{nv_client}{nv_client}{Création / fiche client}%
  \Screenshot{liste_client}{\fichierListeClient}{Liste des clients}%
  \Screenshot{client_dettes}{client_dettes}{Clients et dettes}%
  \Screenshot{demande_devis}{demande_devis}{Demandes de devis}%
  \Screenshot{ajout_demande_devis}{ajout_demande_devis}{Ajout d'une demande de devis}%
  \Screenshot{liste_stock}{\fichierListeStock}{Liste du stock général}%
  \Screenshot{ajout_produit}{ajout_produit}{Ajout de produit}%
  \Screenshot{modif_produit}{modif_produit}{Modification de produit}%
  \Screenshot{supprimer_produit}{supprimer_produit}{Suppression de produit}%
  \Screenshot{facturation_vente}{facturation_vente}{Facturation vente}%
  \Screenshot{nouvelle_facture}{nouvelle_facture}{Création d'une nouvelle facture}%
  \Screenshot{facture_achat}{facture_achat}{Facturation / suivi achat}%
  \Screenshot{fournisseurs}{fourniseurs}{Gestion des fournisseurs (nom de fichier : fourniseurs.png)}%
  \Screenshot{reclamations}{reclamations}{Liste des réclamations}%
  \Screenshot{ajout_reclamation}{ajout_reclamation}{Ajout d'une réclamation}%
  \Screenshot{calendrier}{calendrier}{Calendrier / planning}%
}

% -----------------------------------------------------------------------------
% CHAPITRE 5 — Tests & validation (illustrations des écrans validés)
% Reprendre une sélection représentative (éviter 20 figures d'un coup).
% -----------------------------------------------------------------------------
\newcommand{\FiguresChapitreCinqTests}{%
  \Screenshot{login_test}{login}{Vérification du parcours d'authentification}%
  \Screenshot{dashboard_test}{dashboard}{Contrôle du tableau de bord après connexion}%
  \Screenshot{liste_vehicules_test}{liste_vehicule}{Scénario : navigation liste véhicules}%
  \Screenshot{facturation_test}{facturation_vente}{Scénario : module facturation}%
}
